\section{Immutable Parent Pointers}
\label{sec:immutable-parent-pointers}

Central to the Recursive Spawning mechanism is the \textit{Immutable Parent Pointer} — the structural device that binds every spawned node to its accountable ancestor, prevents any node from operating without a traceable lineage, and makes the upstream accountability guarantee physically enforceable rather than policy-dependent.

\subsection{The ParentPointer Relation}

At the foundation of every spawned BX3 loop is the \textit{ParentPointer} relation, which records the identity of the parent node that authorized and deployed each child. We define this relation formally:

\begin{definition}[ParentPointer]
\label{def:parent-pointer}
Given a parent node $p$ and a child node $c$ produced by $p$'s spawning operation, the relation $\textit{ParentPointer}(p, c)$ is established at the moment of spawn and satisfies two governing properties:

\begin{enumerate}[noitemsep]
    \item \textbf{Uniqueness:} For any child node $c$, there exists \textit{exactly one} parent $p$ such that $\textit{ParentPointer}(p, c)$ holds. No child may have more than one parent, and there is no mechanism by which an existing pointer can be overwritten.
    \item \textbf{Immutability:} Once $\textit{ParentPointer}(p, c)$ is established, it cannot be modified, reassigned, transferred, or severed for the entire operational lifetime of $c$. This holds regardless of what the child node attempts, what any Bounds Engine requests, what any external actor commands, or what hardware or network events occur.
\end{enumerate}
\end{definition}

The uniqueness property is a structural invariant of the spawn operation: when a parent node $p$ generates a child $c$, it records exactly one parent pointer. There is no mechanism by which a second, competing parent pointer can be simultaneously established, and there is no circumstance under which an existing pointer can be overwritten. The child node is born into exactly one lineage.

The immutability property is not a software convention, a permissions policy, or a Bounds Engine constraint. It is a physical property enforced by the Fact Layer, which we examine in full detail in Section~\ref{sec:immutability-enforcement}. For now, we note that immutability is architecturally inseparable from the ParentPointer relation: a pointer that could be changed after instantiation would not satisfy the accountability guarantee, because a node could retroactively sever its lineage and orphan its actions from human accountability.

\subsection{The Accountability Chain}

The ParentPointer relation, applied transitively from any node back to the system origin, generates a structure we call the \textit{accountability chain}. This chain is the complete, ordered ancestry record of any node in the system.

\begin{definition}[Accountability Chain]
\label{def:chain}
For any node $a$ in a spawned BX3 population, the accountability chain of $a$, denoted $\textit{Chain}(a)$, is defined recursively as:

\begin{equation}
\label{eq:chain-recursion}
\textit{Chain}(a) = \textit{ParentPointer}(\textit{parent}(a), a) \;\cup\; \textit{Chain}(\textit{parent}(a))
\end{equation}

\noindent with the base case:

\begin{equation}
\label{eq:chain-base}
\textit{parent}(r) = \bot \quad\text{for the root human node } r \;\Rightarrow\; \textit{Chain}(r) = \emptyset
\end{equation}

\noindent where $\bot$ denotes the null parent, and the root human node $r$ is defined as the single node in the system whose parent is $\bot$.
\end{definition}

Equation~\ref{eq:chain-recursion} states that the chain of any node $a$ consists of the single edge connecting $a$ to its parent, united with the chain of that parent. This is a classic recursive definition with a provably terminating base case: each step moves strictly upward in the hierarchy (from child to parent), and the root has no parent. The result is a finite, ordered set whose elements are the transitive closure of the parent-of relation over the spawn tree.

The chain captures the full ancestry of a node in the spawn hierarchy. Every element in $\textit{Chain}(a)$ is an ancestor of $a$ that participated in authorizing or propagating the Purpose that $a$ currently carries. No actor outside this chain bears accountability for $a$'s current actions, and no actor inside this chain can be removed from it retroactively.

A critical corollary of Definitions \ref{def:parent-pointer} and \ref{def:chain} is the following:

\begin{corrollary}[Unique Human Termination]
Every non-root node $a$ has a non-empty chain that terminates at the root human node $r$. That is, $\forall a \neq r: r \in \textit{Chain}(a)$. The root human is the unique terminal element of every non-empty accountability chain in the system.
\end{corrollary}

This corollary is direct from the recursion: each step moves to $\textit{parent}(a)$, and only the root has $\bot$ as its parent. No spawned node can have an empty chain (it always contains at least its own parent edge), and no chain can terminate anywhere other than the root human.

\subsection{The Root Human Anchor}
\label{sec:root-human}

The root human node is the single human entity at the apex of the accountability hierarchy. In the BX3 Framework, this is the founder, operator, or designated authority who initiates the system and retains ultimate accountability for all downstream spawns. The root human is not a software configuration choice; it is a structural requirement for the system to possess a defined accountability terminus.

The root human node is characterized by the following properties:

\begin{itemize}[noitemsep]
    \item \textbf{Null Parent:} $\textit{parent}(r) = \bot$. The root human is the origin point of all authority and cannot themselves be spawned by any other node within the BX3 system. The root human is root \textit{by definition}: there is nothing above it.
    \item \textbf{Unconditional Accountability:} All chains terminate at the root human. The root bears final accountability for every action taken by every descendant node in the system, whether those actions were taken directly or delegated through intermediate Purpose Layers.
    \item \textbf{Source of Spawn Authorization:} Only the root human --- or a node explicitly and verifiably authorized by the root human --- can authorize the creation of a new spawn event. This authorization is mediated through the Purpose Layer and logged in the Forensic Ledger (Section~\ref{sec:spawn-auth}).
    \item \textbf{Immutable Identity:} The root human node is registered at system initialization and is cryptographically bound to the top-level Purpose Layer. It cannot be replaced, merged, shadowed, or deleted without explicit human action recorded in the Forensic Ledger.
\end{itemize}

In practice, the root human may correspond to a founder-CEO, a board of directors, a regulatory body, or any named institutional authority that can bear legal and operational accountability for the system's actions. The label is flexible; the structural requirement is not. What is required is that a single, unambiguous human anchor exists such that $\forall a: \exists!\, r: r \in \textit{Chain}(a) \cup \{r\}$. Every spawned node in the system traces back to this anchor.

\subsection{Immutability Enforcement by the Fact Layer}
\label{sec:immutability-enforcement}

The immutability of ParentPointer is not a policy, a permissions check, or a Bounds Engine constraint. It is a physical property enforced by the Fact Layer of the node on which the ParentPointer is stored. This distinction is architecturally critical: immutability is structurally guaranteed, not operationally dependent.

Consider the failure modes that immutability must survive:

\begin{itemize}[noitemsep]
    \item \textbf{Bounds Engine compromise:} A Bounds Engine that has been manipulated or is behaving adversarially might attempt to issue a write operation targeting its own or a child node's ParentPointer. The Fact Layer rejects this at the hardware boundary.
    \item \textbf{Child node rebellion:} A spawned child node might attempt to overwrite its own ParentPointer to disavow its lineage. The Fact Layer refuses: the ParentPointer region is read-only at the Fact Layer boundary.
    \item \textbf{Network-level attacks:} Man-in-the-middle attacks, replay attacks, or node impersonation cannot alter an established ParentPointer because the Fact Layer rejects any write operation targeting these values regardless of the requestor's credentials or network position.
    \item \textbf{Hardware failure and recovery:} A node that experiences power loss, hardware damage, or forced reboot does not lose its ParentPointer because the pointer is stored in non-volatile Fact Layer memory. Upon recovery, the node resumes with its chain intact and its immutability guarantees preserved.
    \item \textbf{Administrative override attempts:} Even a human operator with root access at the operating system level cannot modify a Fact-Layer-protected ParentPointer through software means, because the protection is embedded in the execution boundary itself, not at the OS permissions layer.
\end{itemize}

The enforcement model follows directly from the BX3 layer separation:

\begin{enumerate}[noitemsep]
    \item The Fact Layer enforces physical constraints.
    \item ParentPointer immutability is a physical constraint (a region of memory that must not be written after initialization).
    \item Therefore, the Fact Layer enforces ParentPointer immutability as a physical invariant.
\end{enumerate}

This is the same mechanism by which the Fact Layer enforces any hard constraint: it is not a check that can be bypassed by presenting correct credentials, because the enforcement is embedded in the execution boundary itself. The Bounds Engine may model, reason about, simulate, or propose modifications to the chain structure, but it cannot execute any write to the ParentPointer region. The Purpose Layer can authorize new spawn events (creating new ParentPointers), but it cannot edit existing ones. Authorization for spawn and immutability of existing pointers are separate, non-confusable operations that operate at different architectural layers.

\subsection{Spawn Authorization via the Purpose Layer}
\label{sec:spawn-auth}

While existing ParentPointers are immutable, the system must be able to \textit{grow} the accountability chain through new spawn events. The creation of a new ParentPointer is the sole mechanism by which the accountability chain extends. This creation requires explicit, Purpose-Layer-mediated authorization. No node may spontaneously generate a child and thereby establish a new ParentPointer without going through the authorized spawn protocol.

The Purpose Layer mediates spawn authorization through the following canonical sequence:

\begin{enumerate}[noitemsep]
    \item \textbf{Request:} A parent node's Bounds Engine determines that a new child node is required to handle a specific local context or workload. It formulates a spawn request addressed to its own Purpose Layer. This request includes the proposed scope of the child's authority, the local context parameters, and the authorized constraint envelope.
    \item \textbf{Authorization Check:} The Purpose Layer evaluates the request against its registered mission parameters, constraint envelope, and resource allocation policy. It determines whether the parent is authorized to spawn at all, and if so, what scope of Purpose the child will carry. This check is the sole gate for chain extension.
    \item \textbf{Worksheet Generation:} Upon successful authorization, the parent Bounds Engine generates a \textit{Worksheet} --- a containerized, self-contained logic set that encapsulates the authorized Purpose for the specific local context. The Worksheet includes the immutable ParentPointer value (pre-loaded from the parent's Fact Layer) and a cryptographically signed assertion of authorization provenance.
    \item \textbf{Deployment:} The parent deploys the Worksheet to the child node over-the-air (OTA). The child node instantiates its BX3 loop using the received Worksheet, establishing its Fact Layer with the authorized constraint set and the immutable ParentPointer pre-loaded and read-only from the moment of instantiation.
    \item \textbf{Chain Extension:} The new child node is now live with $\textit{ParentPointer}(\text{parent}, \text{child})$ established, immutable, and logged.
\end{enumerate}

A critical architectural property is that the \textit{immutability} of the new ParentPointer begins \textit{at the moment of instantiation}, not at some later checkpoint. From the instant the child node's Fact Layer initializes, the ParentPointer is inextensible and ineditable. This eliminates the possibility of a spawn event occurring in a state where the child has agency but no lineage.

The Purpose Layer's role in spawn authorization ensures that the chain grows in a controlled, auditable manner. Every spawn event is recorded in the Forensic Ledger as a \texttt{CHAIN\_EXTEND} event, capturing the authorizing Purpose Layer reference, the authorized scope transmitted to the child, the child node identifier, the pre-loaded ParentPointer value, and the timestamp.

\subsection{Chain Traversal Algorithm}

The accountability chain must be traversable by any authorized actor in the system, including auditors, regulators, automated compliance systems, and the Purpose Layer itself. We define a canonical chain traversal algorithm that operates with guaranteed termination at the root human node:

\begin{algorithm}[htbp]
\caption{Accountability Chain Traversal}
\label{alg:chain-traversal}
\begin{enumerate}[noitemsep]
    \item \textbf{Input:} A node identifier $a$ for which the full accountability chain is required.
    \item \textbf{Output:} An ordered list $[r, p_1, p_2, \dots, p_n, a]$ where $r$ is the root human node, $p_i$ are intermediate parent nodes, and $a$ is the target node.
    \item \textbf{Procedure:}
    \begin{enumerate}[noitemsep]
        \item Initialize an empty list $\texttt{chain}$.
        \item Set $\texttt{current} \gets a$.
        \item \textbf{while} $\texttt{current} \neq \bot$ \textbf{do}:
        \begin{enumerate}[noitemsep]
            \item $\texttt{append}(\texttt{chain}, \texttt{current})$.
            \item $\texttt{parent} \gets \textit{ReadParent}(\texttt{current})$ \hfill \textit{// Fact Layer read operation}.
            \item \textbf{if} $\texttt{parent} = \bot$ \textbf{then}:
            \begin{enumerate}[noitemsep]
                \item \textbf{assert} $\texttt{current}$ is the root human $r$.
                \item \textbf{break}.
            \end{enumerate}
            \item \textbf{else}: $\texttt{current} \gets \texttt{parent}$.
        \endenumerate}
        \item \textbf{return} $\texttt{reverse}(\texttt{chain})$ \hfill \textit{// root first, target last}.
    \endenumerate}
\end{enumerate}
\end{algorithm}

\textbf{Termination proof.} The while loop terminates because the parent relation is strictly acyclic: by Definition~\ref{def:parent-pointer}, no node can be its own ancestor, and by Definition~\ref{def:chain}, each iteration moves from a node to its parent, which is a strictly weaker position in the spawn DAG. Since the spawn hierarchy has finite depth (bounded by the number of spawn generations), repeated iteration must eventually reach the root, whose parent is $\bot$, triggering the break condition. The loop cannot iterate indefinitely.

\textbf{Correctness proof.} We establish three invariants over the loop:
\begin{itemize}[noitemsep]
    \item \textbf{Invariant 1 (Chain completeness):} At the start of each iteration, $\texttt{chain}$ contains exactly the nodes on the path from $a$ to $\texttt{current}$ (inclusive), in reverse order (from $a$ upward).
    \item \textbf{Invariant 2 (Parent validity):} Every element in $\texttt{chain}$ is a valid descendant of the root human $r$.
    \item \textbf{Invariant 3 (ParentPointer integrity):} The $\textit{ReadParent}()$ operation reads from Fact Layer storage and returns the stored value unchanged; it cannot return a computed or modified value because the Fact Layer enforces immutability at the read boundary.
\end{itemize}
Invariant 1 holds by construction: the algorithm appends $\texttt{current}$ before reading its parent, and each iteration moves to $\texttt{parent}$. Invariant 2 holds because every node in the chain was spawned with a valid ParentPointer, and the chain is closed under the parent operation. Invariant 3 holds by the Fact Layer enforcement described in Section~\ref{sec:immutability-enforcement}. When the loop terminates at $\bot$, reversing $\texttt{chain}$ yields $[r, p_1, p_2, \dots, p_n, a]$, which is exactly the required output.

The reversed ordering --- root first, immediate parent second, target node last --- is the standard representation used in Forensic Ledger entries, regulatory audit reports, and escalation notifications. This ordering makes the chain immediately human-readable: accountability flows from the top down, and any observer can immediately identify the ultimate human anchor for any node's actions.

\subsection{Forensic Ledger Integration}

Every ParentPointer read during chain traversal is logged in the Forensic Ledger as a \texttt{CHAIN\_TRAVERSE} event. The log entry records the node performing the traversal, the target node queried, the ParentPointer values returned at each step, and the timestamp. Because the Fact Layer performs the logging, the entry is itself immutable and deterministic.

Every spawn event that creates a new ParentPointer is logged as a \texttt{CHAIN\_EXTEND} event. This record captures the authorizing Purpose Layer reference, the authorized scope transmitted to the child, the child node identifier, the pre-loaded ParentPointer value, and the timestamp.

Both event types are recorded simultaneously on three discrete planes:

\begin{itemize}[noitemsep]
    \item \textbf{Purpose Plane:} Records the authorization decision, the authorizing actor, and the scope of delegated authority.
    \item \textbf{Bounds Engine Plane:} Records the spawn request, the Bounds Engine's reasoning about the child's proposed context, and the Worksheet contents transmitted.
    \item \textbf{Fact Plane:} Records the physical ParentPointer write operation, the child node's Fact Layer initialization, and the immutable chain state.
\end{itemize}

If any plane's log diverges from the others --- indicating a potential chain integrity violation --- a \texttt{CHAIN\_INTEGRITY\_FAILURE} flag is raised and a Bailout Protocol signal is sent to the nearest Purpose Layer ancestor. This three-plane attestation model ensures that no single plane can be compromised without detection, and no chain integrity event can escape a log record.

The combination of immutability enforcement (Section~\ref{sec:immutability-enforcement}), Purpose Layer spawn authorization (Section~\ref{sec:spawn-auth}), and Fact Layer logged traversal creates an accountability chain that is cryptographically verifiable by construction. Breaking the chain requires simultaneous, undetected compromise of all three functional layers across multiple nodes --- an event with negligibly low probability in any real-world threat model.

\subsection{Summary}

The Immutable Parent Pointers mechanism is the structural guarantee that the BX3 Framework's upstream accountability guarantee is physically enforceable rather than merely policy-dependent. We summarize the core properties:

\begin{itemize}[noitemsep]
    \item \textbf{ParentPointer(p, c)} is established once at spawn and is immutable thereafter. Uniqueness ensures exactly one parent per child. Immutability ensures this relation cannot be corrupted.
    \item \textbf{Chain(a)} is defined recursively as $\textit{ParentPointer}(\textit{parent}(a), a) \cup \textit{Chain}(\textit{parent}(a))$ with base case $\textit{Chain}(r) = \emptyset$ for the root human $r$. The chain is always non-empty for non-root nodes and always terminates at the root human.
    \item \textbf{Immutability enforcement} is a Fact Layer responsibility. The Fact Layer physically cannot execute any write to an established ParentPointer. This is not a permissions check --- it is a hardware or runtime memory boundary that no actor can bypass without structural compromise of the Fact Layer itself.
    \item \textbf{Spawn authorization} is a Purpose Layer responsibility. Only the Purpose Layer can authorize chain extension. Every extension is logged on three planes simultaneously, creating a Forensic Ledger record that is verifiable across all system layers.
    \item \textbf{Chain traversal} is a deterministic, guaranteed-terminating algorithm that recovers the full ancestry of any node. Every traversal reads only from immutable Fact Layer storage and produces an ordered list rooted at the human anchor.
\end{itemize}
